\documentclass[11pt]{article}
% The preamble every project's paper shares: % The preamble every project's paper shares: % The preamble every project's paper shares: \input{../../preamble} from
% projects/<project>/paper/. Boards come from chessboard and xskak, plots from
% pgfplots; every number a paper prints comes from its generated.tex.
\usepackage[LSBC4,T1]{fontenc}
\usepackage[utf8]{inputenc}
\usepackage{lmodern}
\usepackage[margin=1in]{geometry}
\usepackage{microtype}
\usepackage{amsmath}
\usepackage{booktabs}
\usepackage{tabularx}
\usepackage{array}
\usepackage{xcolor}
\usepackage[obeyspaces,spaces]{url}
\usepackage{xskak}
\usepackage{pgfplots}
\pgfplotsset{compat=1.18}
\usepackage{caption}
\usepackage{subcaption}
\usepackage{enumitem}
\usepackage[section]{placeins}
\usepackage[hidelinks]{hyperref}

\setchessboard{boardfontencoding=LSBC4, boardfontfamily=skaknew, blackfieldcolor=black!22,
  whitefieldcolor=white, setfontcolors, showmover=false, label=true, labelfontsize=6pt,
  boardfontsize=14pt}

% A move or line in UCI, and a position in FEN, both breakable at a slash.
\newcommand{\uci}[1]{\texttt{#1}}
\newcommand{\fen}[1]{{\footnotesize\path{#1}}}
\newcolumntype{L}{>{\raggedright\arraybackslash}X}

\setlength{\parskip}{0.4em}
\setlength{\parindent}{0pt}
 from
% projects/<project>/paper/. Boards come from chessboard and xskak, plots from
% pgfplots; every number a paper prints comes from its generated.tex.
\usepackage[LSBC4,T1]{fontenc}
\usepackage[utf8]{inputenc}
\usepackage{lmodern}
\usepackage[margin=1in]{geometry}
\usepackage{microtype}
\usepackage{amsmath}
\usepackage{booktabs}
\usepackage{tabularx}
\usepackage{array}
\usepackage{xcolor}
\usepackage[obeyspaces,spaces]{url}
\usepackage{xskak}
\usepackage{pgfplots}
\pgfplotsset{compat=1.18}
\usepackage{caption}
\usepackage{subcaption}
\usepackage{enumitem}
\usepackage[section]{placeins}
\usepackage[hidelinks]{hyperref}

\setchessboard{boardfontencoding=LSBC4, boardfontfamily=skaknew, blackfieldcolor=black!22,
  whitefieldcolor=white, setfontcolors, showmover=false, label=true, labelfontsize=6pt,
  boardfontsize=14pt}

% A move or line in UCI, and a position in FEN, both breakable at a slash.
\newcommand{\uci}[1]{\texttt{#1}}
\newcommand{\fen}[1]{{\footnotesize\path{#1}}}
\newcolumntype{L}{>{\raggedright\arraybackslash}X}

\setlength{\parskip}{0.4em}
\setlength{\parindent}{0pt}
 from
% projects/<project>/paper/. Boards come from chessboard and xskak, plots from
% pgfplots; every number a paper prints comes from its generated.tex.
\usepackage[LSBC4,T1]{fontenc}
\usepackage[utf8]{inputenc}
\usepackage{lmodern}
\usepackage[margin=1in]{geometry}
\usepackage{microtype}
\usepackage{amsmath}
\usepackage{booktabs}
\usepackage{tabularx}
\usepackage{array}
\usepackage{xcolor}
\usepackage[obeyspaces,spaces]{url}
\usepackage{xskak}
\usepackage{pgfplots}
\pgfplotsset{compat=1.18}
\usepackage{caption}
\usepackage{subcaption}
\usepackage{enumitem}
\usepackage[section]{placeins}
\usepackage[hidelinks]{hyperref}

\setchessboard{boardfontencoding=LSBC4, boardfontfamily=skaknew, blackfieldcolor=black!22,
  whitefieldcolor=white, setfontcolors, showmover=false, label=true, labelfontsize=6pt,
  boardfontsize=14pt}

% A move or line in UCI, and a position in FEN, both breakable at a slash.
\newcommand{\uci}[1]{\texttt{#1}}
\newcommand{\fen}[1]{{\footnotesize\path{#1}}}
\newcolumntype{L}{>{\raggedright\arraybackslash}X}

\setlength{\parskip}{0.4em}
\setlength{\parindent}{0pt}

% Written by projects/tiebreak-rules/paper/generate.mjs; do not edit by hand.

\newcommand{\AppCommit}{ffee8dbb16}
\newcommand{\TrainMonths}{July and August 2026}
\newcommand{\HeldMonth}{September 2026}
\newcommand{\BandLow}{1250}
\newcommand{\HisHigh}{1449}
\newcommand{\AboveLow}{1450}
\newcommand{\BandHigh}{1649}
\newcommand{\Draws}{2,000}
\newcommand{\Seed}{20261021}
\newcommand{\MinHeld}{300}
\newcommand{\MinOrder}{100}
\newcommand{\PerMonth}{30,000}
\newcommand{\Cap}{3,000}
\newcommand{\MinPly}{8}
\newcommand{\MinPlies}{30}
\newcommand{\Nodes}{250,000}
\newcommand{\WindowWp}{3}
\newcommand{\LeadMin}{2}
\newcommand{\LevelWithin}{1}
\newcommand{\TablebaseMen}{seven}
\newcommand{\CoachShare}{80}
\newcommand{\NoInfoShare}{50}
\newcommand{\Games}{90,000}
\newcommand{\Skipped}{0}
\newcommand{\Entering}{2,055,801}
\newcommand{\InUnit}{837,954}
\newcommand{\MaiaCalls}{837,954}
\newcommand{\Scored}{54,353}
\newcommand{\UnitRows}{59,616}
\newcommand{\GamesTrain}{24,121}
\newcommand{\GamesHeld}{12,043}
\newcommand{\ThinName}{ahead in development}
\newcommand{\ThinHeld}{1,196}
\newcommand{\ThinHeldFollow}{329}
\newcommand{\CoverageTable}{%
\begin{table}[ht]
\centering\footnotesize
\begin{tabular}{@{}lrrrrrr@{}}
\toprule
 & \begin{tabular}[b]{@{}l@{}}In situation, \\ held out\end{tabular} & \begin{tabular}[b]{@{}l@{}}Split, \\ held out\end{tabular} & \begin{tabular}[b]{@{}l@{}}Kept, \\ train\end{tabular} & \begin{tabular}[b]{@{}l@{}}Kept, \\ held out\end{tabular} & \begin{tabular}[b]{@{}l@{}}Follow- \\ first\end{tabular} & \begin{tabular}[b]{@{}l@{}}Break- \\ first\end{tabular} \\
\midrule
Tension & 49,779 & 22,180 & 6,000 & 3,000 & 960 & 2,040 \\
Ahead in development & 3,359 & 1,454 & 2,462 & 1,196 & 329 & 867 \\
Behind in development & 6,651 & 2,888 & 4,879 & 2,383 & 1,565 & 818 \\
Bishop for knight & 48,681 & 13,805 & 6,000 & 3,000 & 1,518 & 1,482 \\
Knight for bishop & 8,975 & 6,306 & 6,000 & 3,000 & 2,308 & 692 \\
Bishop's colour & 189,254 & 7,644 & 6,000 & 3,000 & 1,501 & 1,499 \\
\midrule
Level development & 60,057 & 28,609 & 6,000 & 3,000 & 1,012 & 1,988 \\
The control & 267,036 & 1,442 & 2,436 & 1,260 & 799 & 461 \\
\bottomrule
\end{tabular}
\caption{Where each unit comes up. In situation: the held-out quiet positions in the unit's situation. Split: those where one of the two natural moves follows the unit's rule and the other breaks it. Kept: at most one position a game, drawn at random among its splits, and at most 3,000 a month. The last two columns divide the held-out kept positions by the move Maia-2 rates first: the one that follows, or the one that breaks.}
\label{tab:coverage}
\end{table}
}
\newcommand{\ControlTrain}{29.1}
\newcommand{\ControlTrainCi}{26.3 to 31.8}
\newcommand{\ControlTrainN}{2,436}
\newcommand{\ControlOneTrain}{76}
\newcommand{\ControlHeld}{28.5}
\newcommand{\ControlHeldCi}{24.4 to 32.4}
\newcommand{\ControlHeldN}{1,260}
\newcommand{\ControlOneHeld}{75}
\newcommand{\TenTrain}{$-$0.4}
\newcommand{\TenTrainCi}{$-$2.2 to 1.5}
\newcommand{\TenHeld}{$-$1.5}
\newcommand{\TenHeldCi}{$-$4.1 to 1.1}
\newcommand{\TrTrain}{4.5}
\newcommand{\TrTrainCi}{2.6 to 6.3}
\newcommand{\TrHeld}{1.1}
\newcommand{\TrHeldCi}{$-$1.8 to 3.9}
\newcommand{\BkTrain}{5.6}
\newcommand{\BkTrainCi}{4.3 to 6.9}
\newcommand{\BkHeld}{4.0}
\newcommand{\BkHeldCi}{2.1 to 5.9}
\newcommand{\ColTrain}{3.4}
\newcommand{\ColTrainCi}{1.8 to 5.0}
\newcommand{\ColHeld}{4.0}
\newcommand{\ColHeldCi}{1.7 to 6.1}
\newcommand{\BkBfHeld}{17}
\newcommand{\ColBfHeld}{10}
\newcommand{\OneLow}{54}
\newcommand{\OneHigh}{55}
\newcommand{\BkOneTrain}{55}
\newcommand{\BkOneHeld}{54}
\newcommand{\ColOneTrain}{54}
\newcommand{\ColOneHeld}{54}
\newcommand{\AgainstLow}{5}
\newcommand{\AgainstHigh}{12}
\newcommand{\RuleTable}{%
\begin{table}[ht]
\centering\footnotesize
\setlength{\tabcolsep}{4pt}
\begin{tabularx}{\textwidth}{@{}llllL@{}}
\toprule
Rule & \begin{tabular}[b]{@{}l@{}}Balanced lift, \\ train\end{tabular} & \begin{tabular}[b]{@{}l@{}}Balanced lift, \\ held out\end{tabular} & \begin{tabular}[b]{@{}l@{}}Break-first lift, \\ held out\end{tabular} & Verdict \\
\midrule
Keep the tension & $-$0.4 {\scriptsize($-$2.2, 1.5)} & $-$1.5 {\scriptsize($-$4.1, 1.1)} & $-$17.2 {\scriptsize($-$20.0, $-$14.3)} & does no better than Maia-2's order \\
Trades and development & 4.5 {\scriptsize(2.6, 6.3)} & 1.1 {\scriptsize($-$1.8, 3.9)} & $-$16.1 {\scriptsize($-$19.4, $-$12.7)} & does no better than Maia-2's order \\
Bishop over knight & 5.6 {\scriptsize(4.3, 6.9)} & 4.0 {\scriptsize(2.1, 5.9)} & $-$16.6 {\scriptsize($-$19.6, $-$13.5)} & breaks the tie \\
Pawns off the bishop's colour & 3.4 {\scriptsize(1.8, 5.0)} & 4.0 {\scriptsize(1.7, 6.1)} & $-$10.3 {\scriptsize($-$13.4, $-$7.0)} & breaks the tie \\
\midrule
The control & 29.1 {\scriptsize(26.3, 31.8)} & 28.5 {\scriptsize(24.4, 32.4)} & 16.3 {\scriptsize(9.5, 22.9)} & passes \\
\bottomrule
\end{tabularx}
\caption{The verdicts. The balanced lift is the mean, over the two orders, of the share of positions where the move that follows the rule is sound less the share where the move that breaks it is, in points, with its 95\% interval from 2,000 draws over games; a rule of two units takes the mean of its units' lifts on every draw. The break-first lift is the same difference where Maia-2 rates the move that breaks the rule first: a rule is worth switching to when its interval lies above 0 on both periods, and none does.}
\label{tab:rules}
\end{table}
}
\newcommand{\BfnTrain}{9.5}
\newcommand{\BfnTrainCi}{7.7 to 11.1}
\newcommand{\BfnHeld}{8.1}
\newcommand{\BfnHeldCi}{5.8 to 10.7}
\newcommand{\NfbTrain}{1.7}
\newcommand{\NfbTrainCi}{$-$0.4 to 3.8}
\newcommand{\NfbHeld}{$-$0.1}
\newcommand{\NfbHeldCi}{$-$3.1 to 2.9}
\newcommand{\BehindTrain}{5.1}
\newcommand{\BehindTrainCi}{3.0 to 7.0}
\newcommand{\BehindHeld}{3.8}
\newcommand{\BehindHeldCi}{1.0 to 6.6}
\newcommand{\AheadTrain}{3.9}
\newcommand{\AheadTrainCi}{0.8 to 6.9}
\newcommand{\AheadHeld}{$-$1.5}
\newcommand{\AheadHeldCi}{$-$6.4 to 3.1}
\newcommand{\LevelTrain}{0.7}
\newcommand{\LevelTrainCi}{$-$1.1 to 2.5}
\newcommand{\LevelHeld}{$-$1.1}
\newcommand{\LevelHeldCi}{$-$3.7 to 1.5}
\newcommand{\UnitTable}{%
\begin{table}[ht]
\centering\scriptsize
\setlength{\tabcolsep}{3pt}
\begin{tabular}{@{}lrrrlllr@{}}
\toprule
 & \begin{tabular}[b]{@{}l@{}}Held-out \\ positions\end{tabular} & \begin{tabular}[b]{@{}l@{}}F \\ sound\end{tabular} & \begin{tabular}[b]{@{}l@{}}B \\ sound\end{tabular} & \begin{tabular}[b]{@{}l@{}}Lift, \\ follow-first\end{tabular} & \begin{tabular}[b]{@{}l@{}}Lift, \\ break-first\end{tabular} & \begin{tabular}[b]{@{}l@{}}Balanced \\ lift\end{tabular} & \begin{tabular}[b]{@{}l@{}}F when only \\ one is sound\end{tabular} \\
\midrule
Tension & 3,000 & 57 & 65 & 14.1 {\scriptsize(9.7, 18.3)} & $-$17.2 {\scriptsize($-$20.0, $-$14.3)} & $-$1.5 {\scriptsize($-$4.1, 1.1)} & 49 \\
Ahead in development & 1,196 & 52 & 59 & 10.9 {\scriptsize(2.9, 19.4)} & $-$14.0 {\scriptsize($-$18.6, $-$9.1)} & $-$1.5 {\scriptsize($-$6.4, 3.1)} & 48 \\
Behind in development & 2,383 & 67 & 56 & 25.8 {\scriptsize(22.5, 29.1)} & $-$18.2 {\scriptsize($-$23.0, $-$13.6)} & 3.8 {\scriptsize(1.0, 6.6)} & 53 \\
Bishop for knight & 3,000 & 63 & 54 & 27.1 {\scriptsize(23.6, 30.5)} & $-$10.8 {\scriptsize($-$14.2, $-$7.1)} & 8.1 {\scriptsize(5.8, 10.7)} & 57 \\
Knight for bishop & 3,000 & 66 & 54 & 22.2 {\scriptsize(19.3, 25.0)} & $-$22.4 {\scriptsize($-$27.7, $-$17.3)} & $-$0.1 {\scriptsize($-$3.1, 2.9)} & 50 \\
Bishop's colour & 3,000 & 67 & 63 & 18.3 {\scriptsize(15.0, 21.5)} & $-$10.3 {\scriptsize($-$13.4, $-$7.0)} & 4.0 {\scriptsize(1.7, 6.1)} & 54 \\
\midrule
Level development & 3,000 & 55 & 63 & 19.2 {\scriptsize(15.0, 23.4)} & $-$21.4 {\scriptsize($-$24.6, $-$18.3)} & $-$1.1 {\scriptsize($-$3.7, 1.5)} & 49 \\
The control & 1,260 & 62 & 30 & 40.7 {\scriptsize(36.0, 45.0)} & 16.3 {\scriptsize(9.5, 22.9)} & 28.5 {\scriptsize(24.4, 32.4)} & 75 \\
\bottomrule
\end{tabular}
\caption{Each unit apart, held out. F and B sound: the share of positions where the move that follows and the move that breaks the unit's rule is sound, in percent. The lifts are F's share less B's in points, in each order and as their mean, with 95\% intervals. F when only one is sound: among the positions where exactly one of the two is sound, the share where it is F, as the mean of the two orders' shares, in percent.}
\label{tab:units}
\end{table}
}
\newcommand{\BfnCloseBfTrain}{4.2}
\newcommand{\BfnCloseBfTrainCi}{0.8 to 7.6}
\newcommand{\BfnCloseBfHeld}{2.4}
\newcommand{\BfnCloseBfHeldCi}{$-$2.2 to 7.3}
\newcommand{\BfnClearBfHeld}{26}
\newcommand{\ColCloseBfTrain}{$-$1.5}
\newcommand{\ColCloseBfTrainCi}{$-$4.6 to 1.4}
\newcommand{\ColCloseBfHeld}{$-$2.1}
\newcommand{\ColCloseBfHeldCi}{$-$6.2 to 2.2}
\newcommand{\ColClearBfHeld}{20}
\newcommand{\LevelCloseTrain}{4.9}
\newcommand{\LevelCloseHeld}{3.8}
\newcommand{\TenCloseTrain}{4.3}
\newcommand{\TenCloseHeld}{3.2}
\newcommand{\CloseShareLow}{41}
\newcommand{\CloseShareHigh}{56}
\newcommand{\BfnOpenHeld}{10.2}
\newcommand{\BfnMiddleHeld}{7.0}
\newcommand{\ColOpenHeld}{$-$2.7}
\newcommand{\ColMiddleHeld}{5.5}
\newcommand{\ColEndHeld}{6.0}
\newcommand{\NfbOpenHeld}{8.5}
\newcommand{\NfbMiddleHeld}{$-$4.4}
\newcommand{\BfnHisHeld}{11.5}
\newcommand{\BfnAboveHeld}{4.9}
\newcommand{\BehindBishopHeld}{21.5}
\newcommand{\BehindKnightHeld}{8.6}
\newcommand{\CloseTable}{%
\begin{table}[ht]
\centering\scriptsize
\setlength{\tabcolsep}{2pt}
\begin{tabular}{@{}lrllrll@{}}
\toprule
 & \multicolumn{3}{c}{Close calls} & \multicolumn{3}{c}{Clear first choice} \\
\cmidrule(lr){2-4}\cmidrule(l){5-7}
 & Positions & \begin{tabular}[b]{@{}l@{}}Balanced \\ lift\end{tabular} & \begin{tabular}[b]{@{}l@{}}Break-first \\ lift\end{tabular} & Positions & \begin{tabular}[b]{@{}l@{}}Balanced \\ lift\end{tabular} & \begin{tabular}[b]{@{}l@{}}Break-first \\ lift\end{tabular} \\
\midrule
Tension & 1,414 & 3.2 {\scriptsize($-$0.4, 6.8)} & $-$2.8 {\scriptsize($-$7.3, 1.7)} & 1,586 & $-$2.5 {\scriptsize($-$6.3, 1.2)} & $-$26.9 {\scriptsize($-$30.6, $-$23.0)} \\
Ahead in development & 563 & 2.1 {\scriptsize($-$4.4, 8.2)} & $-$0.3 {\scriptsize($-$7.8, 7.4)} & 633 & 0.5 {\scriptsize($-$7.0, 8.0)} & $-$22.9 {\scriptsize($-$29.2, $-$16.3)} \\
Behind in development & 1,167 & $-$1.8 {\scriptsize($-$6.0, 2.2)} & $-$12.6 {\scriptsize($-$18.6, $-$6.7)} & 1,216 & 5.0 {\scriptsize(0.4, 9.3)} & $-$27.8 {\scriptsize($-$35.8, $-$19.9)} \\
Bishop for knight & 1,673 & 9.8 {\scriptsize(6.4, 13.1)} & 2.4 {\scriptsize($-$2.2, 7.3)} & 1,327 & 7.2 {\scriptsize(3.7, 10.8)} & $-$26.2 {\scriptsize($-$31.0, $-$21.1)} \\
Knight for bishop & 1,219 & $-$6.2 {\scriptsize($-$10.3, $-$2.2)} & $-$17.4 {\scriptsize($-$23.9, $-$11.0)} & 1,781 & $-$1.7 {\scriptsize($-$6.5, 3.2)} & $-$33.6 {\scriptsize($-$42.7, $-$24.8)} \\
Bishop's colour & 1,659 & 3.4 {\scriptsize(0.3, 6.5)} & $-$2.1 {\scriptsize($-$6.2, 2.2)} & 1,341 & 5.2 {\scriptsize(1.9, 8.5)} & $-$19.9 {\scriptsize($-$24.8, $-$14.9)} \\
\midrule
Level development & 1,409 & 3.8 {\scriptsize(0.2, 7.5)} & $-$3.0 {\scriptsize($-$7.9, 2.1)} & 1,591 & 0.7 {\scriptsize($-$3.1, 4.6)} & $-$32.8 {\scriptsize($-$36.8, $-$29.1)} \\
The control & 718 & 26.5 {\scriptsize(21.4, 31.6)} & 21.4 {\scriptsize(13.8, 28.7)} & 542 & 26.1 {\scriptsize(19.0, 33.0)} & 3.1 {\scriptsize($-$10.0, 16.2)} \\
\bottomrule
\end{tabular}
\caption{Close calls against a clear first choice, held out, in points with 95\% intervals. A close call: Maia-2's probability of its second move is at least half its probability of its first. The break-first lift is F's soundness less B's where Maia-2 rates B first. The splits describe and decide nothing.}
\label{tab:close}
\end{table}
}
\newcommand{\SplitTable}{%
\begin{table}[ht]
\centering\scriptsize
\setlength{\tabcolsep}{4pt}
\begin{tabular}{@{}lrrrrr@{}}
\toprule
 & \begin{tabular}[b]{@{}l@{}}1250 to \\ 1449\end{tabular} & \begin{tabular}[b]{@{}l@{}}1450 to \\ 1649\end{tabular} & Opening & \begin{tabular}[b]{@{}l@{}}Middle- \\ game\end{tabular} & Endgame \\
\midrule
Tension & $-$2 {\scriptsize(1,299)} & $-$1 {\scriptsize(1,701)} & $-$2 {\scriptsize(1,304)} & $-$3 {\scriptsize(1,261)} & 5 {\scriptsize(435)} \\
Ahead in development & $-$1 {\scriptsize(538)} & $-$2 {\scriptsize(658)} & $-$4 {\scriptsize(763)} & 2 {\scriptsize(433)} & -- \\
Behind in development & 0 {\scriptsize(1,083)} & 7 {\scriptsize(1,300)} & 2 {\scriptsize(1,598)} & 7 {\scriptsize(785)} & -- \\
Bishop for knight & 12 {\scriptsize(1,450)} & 5 {\scriptsize(1,550)} & 10 {\scriptsize(1,354)} & 7 {\scriptsize(1,509)} & $-$3 {\scriptsize(137)} \\
Knight for bishop & $-$2 {\scriptsize(1,275)} & 1 {\scriptsize(1,725)} & 8 {\scriptsize(1,061)} & $-$4 {\scriptsize(1,715)} & $-$6 {\scriptsize(224)} \\
Bishop's colour & 4 {\scriptsize(1,364)} & 4 {\scriptsize(1,636)} & $-$3 {\scriptsize(572)} & 6 {\scriptsize(1,704)} & 6 {\scriptsize(724)} \\
\midrule
Level development & 1 {\scriptsize(1,332)} & $-$3 {\scriptsize(1,668)} & 0 {\scriptsize(1,450)} & $-$3 {\scriptsize(1,550)} & -- \\
The control & 30 {\scriptsize(572)} & 27 {\scriptsize(688)} & 23 {\scriptsize(211)} & 30 {\scriptsize(930)} & 24 {\scriptsize(119)} \\
\bottomrule
\end{tabular}
\caption{The balanced lift by the mover's band and by phase, held out, in points, with the cell's positions in brackets; a dash is a cell with none. The bands are the mover's Lichess rating; the phase is the app's divider's. The splits describe and decide nothing; every cell's interval is in \texttt{results/rules.json}.}
\label{tab:splits}
\end{table}
}
\newcommand{\BfnExamplesF}{three}
\newcommand{\BfnExamplesB}{three}
\newcommand{\ColExamplesF}{four}
\newcommand{\ColExamplesB}{two}
\newcommand{\BfnFourSide}{White}
\newcommand{\BfnFourBand}{the band above his}
\newcommand{\BfnFourF}{6.~e3}
\newcommand{\BfnFourB}{6.~\symbishop{}xe7}
\newcommand{\BfnFourPlayed}{6.~c4}
\newcommand{\BfnFourBest}{6.~\symknight{}f3}
\newcommand{\BfnFourFWin}{43}
\newcommand{\BfnFourBWin}{35}
\newcommand{\ColTwoSide}{Black}
\newcommand{\ColTwoBand}{his band}
\newcommand{\ColTwoF}{14\ldots{}b6}
\newcommand{\ColTwoB}{14\ldots{}d5}
\newcommand{\ColTwoPlayed}{14\ldots{}b6}
\newcommand{\ColTwoBest}{14\ldots{}b6}
\newcommand{\ColTwoFWin}{69}
\newcommand{\ColTwoBWin}{64}
\newcommand{\BfnSixSide}{Black}
\newcommand{\BfnSixBand}{the band above his}
\newcommand{\BfnSixF}{24\ldots{}\symrook{}c8}
\newcommand{\BfnSixB}{24\ldots{}\symbishop{}xf5}
\newcommand{\BfnSixPlayed}{24\ldots{}\symrook{}c8}
\newcommand{\BfnSixBest}{24\ldots{}\symbishop{}xf5}
\newcommand{\BfnSixFWin}{24}
\newcommand{\BfnSixBWin}{41}
\newcommand{\ExampleFigure}{%
\begin{figure}[ht]
\centering
\begin{subfigure}[t]{0.32\textwidth}
\centering
\chessboard[setfen={rnbq1rk1/ppppnppp/1b6/1P2p3/Q7/B1P5/P2PPPPP/RN2KBNR w KQ - 1 6}, boardfontsize=11pt, pgfstyle=straightmove, arrow=to, linewidth=0.07em, color=red!75!black, markmoves={a3-e7}, pgfstyle=straightmove, arrow=to, linewidth=0.07em, color=green!50!black, markmoves={e2-e3}]
\caption{Bishop for knight, break-first: 6.~e3 is sound, 6.~\symbishop{}xe7 is not.}
\end{subfigure}\hfill
\begin{subfigure}[t]{0.32\textwidth}
\centering
\chessboard[setfen={r1b1nrk1/pp1p1pp1/2q1p2p/2p5/2P1N3/2PP3P/P4PP1/RN1Q1RK1 b - - 2 14}, boardfontsize=11pt, pgfstyle=straightmove, arrow=to, linewidth=0.07em, color=red!75!black, markmoves={d7-d5}, pgfstyle=straightmove, arrow=to, linewidth=0.07em, color=green!50!black, markmoves={b7-b6}]
\caption{The bishop's colour, break-first: 14\ldots{}b6 is sound, 14\ldots{}d5 is not.}
\end{subfigure}\hfill
\begin{subfigure}[t]{0.32\textwidth}
\centering
\chessboard[setfen={1k1r3r/1p1q1p2/1Q1p1nb1/1n2pN2/pP2P1P1/P2P3P/1B2BP2/R3K2R b KQ - 1 24}, boardfontsize=11pt, pgfstyle=straightmove, arrow=to, linewidth=0.07em, color=red!75!black, markmoves={d8-c8}, pgfstyle=straightmove, arrow=to, linewidth=0.07em, color=green!50!black, markmoves={g6-f5}]
\caption{Bishop for knight, follow-first: 24\ldots{}\symbishop{}xf5 is sound, 24\ldots{}\symrook{}c8 is not.}
\end{subfigure}
\caption{Three of the held-out examples, drawn at random among the positions where exactly one of the two natural moves is sound. Green: the sound move; red: the unsound one.}
\label{fig:examples}
\end{figure}
}


\title{Does a coach's rule pick the better of his level's two natural moves?\\[0.3em]
\large Four rules for the quiet moments of a game, as tie-breakers between the two moves Maia-2
rates most likely, in \Games{} Lichess games refereed by Stockfish}
\author{Chess Coach research lab}
\date{}

\begin{document}
\maketitle

\begin{abstract}
A coach's rules for the quiet moments of a game are not meant to choose a move alone: his student
names a few normal moves, and the rule says which of them to play. This study asks four such rules,
from a titled player's review of a club player's rapid games, only to break that tie. In quiet
positions from \Games{} Lichess games in the player's band and the one above, where the two moves
Maia-2 rates most likely at the mover's band are one that follows a rule and one that breaks it, is
the one that follows sound more often than Maia-2's own order accounts for? The rule, fixed before
any position was scored, first asks a planted control to show that the design sees a rule known to
hold, and it does: a move that leaves a piece to be taken for a gain loses the tie by
\ControlTrain{} points on the train months and \ControlHeld{} held out. Two rules then break the
tie, their intervals above 0 on both periods: bishop over knight, by \BkTrain{} and \BkHeld{}
points, and pawns off the bishop's colour, by \ColTrain{} and \ColHeld. Keeping the tension does no
better than Maia-2's order, and trading by development breaks the tie on the train months only.
Neither winner overrules a first choice. Where Maia-2 rates the move that breaks the rule first,
that move is still sound more often than the rule's, by \BkBfHeld{} and \ColBfHeld{} points held
out, and where only one of the two is sound, bishop over knight's move is that one \BkOneHeld\% of
the time held out, not the coach's four games in five. The two rules choose between two moves that
look about equally good; they are no reason to play the one that looks worse.
\end{abstract}

\section{The question}

A titled player reviewing the rapid games of an adult player rated about 1100 gives rules for the
quiet moments of a game, when nothing is hanging and no tactic is on the board: keep the tension
unless taking gains something; avoid even trades when ahead in development and look for them when
behind; keep the pawns off the bishop's colour; and prefer the bishop to the knight, the better
trade, he says, about four games in five.\footnote{\url{https://www.youtube.com/watch?v=PcGDehRSqTM}}
His fifth rule, the central break while the other king is still in the centre, was tested in an
earlier study and showed no signal.

Two earlier studies asked a rule of thumb to choose the move: among the legal moves that show the
rule, the one Maia-2 rates most likely at the mover's band. The rule's move, one his level plays
less often, was sound less often than the move his level would naturally play, before the tactics
were checked and after. The coach does not use his rules that way. The student names a few normal
moves, and the rule says which of them to play. So this study asks a rule only to break a tie. In
quiet positions where the two moves Maia-2 rates most likely at his level are one that follows the
rule and one that breaks it, is the one that follows it sound more often, beyond what Maia-2's own
order says?

If a rule breaks the tie, a lesson can teach it as the way to choose between two natural moves, and
Game Review can name it where he played the move that breaks it and his other natural move followed
it. If it does no better than Maia-2's order, it explains a move but does not choose between his
candidates. If it misleads, a lesson that teaches it should say where it fails.

\section{How it was measured}

\paragraph{The rules.} Four rules, read as six units (Table~\ref{tab:design}). A unit has a
situation and reads each legal move: the move \emph{follows} the unit's rule, \emph{breaks} it, or
neither. A rule of two units reads two mirrored halves, so that a liking for quiet moves over
trades in general, or for trades over quiet moves, cancels between them. Every reading is the app's
knowledge service at commit \texttt{\AppCommit}. A \emph{quiet move} is not a capture and gives no
check. An \emph{even trade} is a capture that gives no check and that the app's capture war scores
0: what it takes, less the opponent's best continuation on its square. A \emph{lever} is a pawn of
the mover's that attacks a pawn of the opponent's. Development has no reading in the app, so it is
defined here: a side's undeveloped count is its knights and bishops on their starting squares, plus
one when its king stands on its starting square, and the mover's lead is the opponent's count less
the mover's. Material is level when the mover's points are within \LevelWithin{} of the
opponent's, which keeps out the rules about trading when ahead or behind in material.

\begin{table}[ht]
\centering\footnotesize
\begin{tabularx}{\textwidth}{@{}lLLL@{}}
\toprule
Unit & Situation & Follows & Breaks \\
\midrule
Tension & the mover has a lever whose capture is an even trade & a quiet move that moves no lever's pawn & a lever's capture that is an even trade \\
Ahead in development & not an endgame, material level, an even trade of a piece on, the mover \LeadMin{} or more ahead in development & a quiet move & an even trade of a piece \\
Behind in development & the same, the mover \LeadMin{} or more behind & an even trade of a piece & a quiet move \\
Bishop for knight & the mover can give a bishop for a knight in an even trade, and cannot take a bishop with a knight so & a quiet move & the bishop taking the knight \\
Knight for bishop & the mover can take a bishop with a knight in an even trade, and cannot give a bishop for a knight so & the knight taking the bishop & a quiet move \\
Bishop's colour & the mover has one bishop, and quiet pawn moves to squares of both colours & a quiet pawn move to a square of the other colour & a quiet pawn move to a square of the bishop's colour \\
\midrule
Level development & as ahead in development, with development level & a quiet move & an even trade of a piece \\
The control & any position where a rule's unit is in its situation & a quiet move after which the opponent has no winning capture & a quiet piece move that leaves the piece to be taken for a gain \\
\bottomrule
\end{tabularx}
\caption{The units. The six above the line read the four rules: keep the tension and pawns off
the bishop's colour are a unit each, trades and development reads its two halves, ahead and behind
in development, and bishop over knight its two, bishop for knight and knight for bishop. Level
development is a reference, and no rule: it shows how a quiet move fares against an even trade when
development says nothing. The control is planted: a move that leaves its piece to be taken for a
gain is known to be the worse of two far more often than not.}
\label{tab:design}
\end{table}

\paragraph{The positions.} Lichess rated games at 10 minutes with no increment, his own time
control, that ended normally or on time, with at least \MinPlies{} plies, both players rated
\BandLow{} to \BandHigh{}: the first \PerMonth{} such games of each month, \TrainMonths{} to train
on and \HeldMonth{} held out. A position enters from ply \MinPly{} on when the move just played
captured nothing, the mover is not in check and has two legal moves or more, and the app gives
neither side a mate in one or a winning capture. $M_1$ and $M_2$ are the two moves Maia-2, the
app's, rates most likely at the mover's band, \BandLow{} to \HisHigh{} (his) or the band above. A
position \emph{splits} for a unit when it is in the unit's situation and one of $M_1$ and $M_2$
follows the unit's rule and the other breaks it: $F$ is the move that follows and $B$ the move that
breaks. It is \emph{follow-first} when $M_1$ is $F$ and \emph{break-first} when $M_1$ is $B$. Each
game gives a unit at most one position, drawn at random among its splits, and each unit keeps the
first \Cap{} a month.

\paragraph{The referee.} With \TablebaseMen{} men or fewer the app's tablebase probe decides, and a
move is \emph{sound} when it keeps the mover's outcome. Otherwise Stockfish, through the app's
engine service at its grading budget of \Nodes{} nodes with the hash cleared for each position,
searches the best and then scores $F$ and $B$, and a move is sound when it gives up no more than
\WindowWp{} win\% against that best. The move played is recorded and not scored.

\paragraph{The run.} The rule, every script, the app's code and its installed packages were frozen
by hash before any position was scored, and each stage of the run checked the freeze before it
started. Once, between two stages, the check found cache files that another checkout's test runner,
sharing the run's installed packages on the same machine, had written among them. Only that test
runner reads them; once they were moved aside the freeze held again, and the run resumed from its
checkpoints.

\paragraph{The measures.} The \emph{lift} is $F$'s share sound less $B$'s, paired by position, in
points. The move Maia-2 rates first is likely to be sound more often than its second, which adds
to the lift in the follow-first positions and takes from it in the break-first ones. The
\emph{balanced lift} is the mean of the two orders' lifts, which cancels that advantage when it is
the same in both orders, so a rule that carries no information reads about 0 whatever share of
its splits come in each order. A rule's lift is its unit's balanced lift, or the mean of its two
units' on every draw. \emph{When only one is sound} is the share of the positions with exactly one
of $F$ and $B$ sound where it is $F$, the mean of the two orders' shares: a rule with no
information reads \NoInfoShare\%, and the coach's four games in five would read \CoachShare\%.
\emph{Against the first choice} is $F$'s share sound less $M_1$'s: what playing the rule's move in
every split would gain over playing Maia-2's first choice.

\paragraph{The rule, fixed before any position was scored.} Intervals are 95\%, from \Draws{}
bootstrap draws over games (seed \Seed), within each period, every unit's measures on a draw using
the positions of the games drawn. The control comes first: its balanced lift must have its interval
above 0 on both periods, with enough data, or no rule receives a verdict. Then a rule \emph{breaks
the tie} when its lift's interval lies above 0 on the train months and on the held-out month,
\emph{misleads} when it lies below 0 on both, and \emph{does no better than Maia-2's order}
otherwise. A rule whose units have fewer than \MinHeld{} held-out positions, or fewer than
\MinOrder{} in either order, has not enough data. Beside the verdict, and changing none, a rule is
\emph{worth switching to} when its lift in the break-first positions has its interval above 0 on
both periods: there his level's first move breaks the rule and its second follows it, so the rule
is a reason to play the second.

\section{What was scored}

The scan read the first \PerMonth{} games of each month, \Games{} in all. \Entering{} positions
entered it, and Maia-2 read the \InUnit{} of them that stand in some unit's situation, and no
other. The draw and the cap kept \Scored{} positions, from \GamesTrain{} games on the train months
and \GamesHeld{} held out; a position can serve several units, so the units read \UnitRows{} in
all. Tension, bishop for knight, knight for bishop and the bishop's colour filled the cap of
\Cap{} every month (Table~\ref{tab:coverage}). The rarest split is being \ThinName: \ThinHeld{}
positions held out, \ThinHeldFollow{} of them follow-first. So every unit passes the rule's counts
of \MinHeld{} held-out positions and \MinOrder{} in each order.

\CoverageTable

\section{The control passes}

The control's balanced lift is \ControlTrain{} points (\ControlTrainCi) over \ControlTrainN{}
positions on the train months and \ControlHeld{} (\ControlHeldCi) over \ControlHeldN{} held out,
both intervals above 0: the design sees a rule known to hold, and the rules receive verdicts. Where
only one of its two moves is sound, it is the move that leaves nothing to be taken
\ControlOneTrain\% of the time on the train months and \ControlOneHeld\% held out. The control
holds against Maia-2's order as well: where Maia-2 rates the move that leaves a piece to be taken
first, its second move is sound more often (Table~\ref{tab:rules}). Every rule's interval lies below
the control's on both periods, so what follows is about rules that tip a choice, not rules that
decide one.

\section{Two rules break the tie}

Bishop over knight breaks the tie by \BkTrain{} points (\BkTrainCi) on the train months and
\BkHeld{} (\BkHeldCi) held out, and pawns off the bishop's colour by \ColTrain{} (\ColTrainCi) and
\ColHeld{} (\ColHeldCi) (Table~\ref{tab:rules}). Keeping the tension does no better than Maia-2's
order, at \TenTrain{} points (\TenTrainCi) and \TenHeld{} (\TenHeldCi). Trading by development
breaks the tie on the train months, by \TrTrain{} points (\TrTrainCi), but not held out, at
\TrHeld{} (\TrHeldCi); the periods disagree, so by the rule it does no better than Maia-2's order
either.

\RuleTable

\section{A tie-breaker, not a reason to switch}

The balanced lift is half the difference between two leads: how far Maia-2's first choice leads
its second in soundness where the first follows the rule, and how far it leads where the first
breaks it. For the two winners the first choice leads by more where it follows the rule, which is
the rule's information, but it leads in both orders, on both periods (Table~\ref{tab:units}). So
neither winner is worth switching to, and no rule is: every rule's break-first lift lies below 0 on
both periods. Where Maia-2 rates the move that breaks the rule first, that move is sound more often
than the second, the rule's, by \BkBfHeld{} points for bishop over knight and \ColBfHeld{} for the
bishop's colour held out.

The other measures say the same. Where only one of the two moves is sound, it is the rule's move
\BkOneTrain\% of the time on the train months and \BkOneHeld\% held out for bishop over knight, and
\ColOneTrain\% and \ColOneHeld\% for the bishop's colour: above the \NoInfoShare\% of a rule with
no information, and below the \CoachShare\% that the coach's four games in five would read. And
playing the rule's move in every split, in place of Maia-2's first choice, costs \AgainstLow{} to
\AgainstHigh{} points of soundness held out across the four rules, every interval below 0 on both
periods.

\section{The halves}

A rule of two units is read as the mean of its halves, because a lean towards quiet moves, or
towards trades, lifts one half and sinks the other by as much. The halves are reported apart and
decide nothing (Table~\ref{tab:units}).

Bishop over knight's halves differ. Keeping the bishop rather than giving it for a knight lifts by
\BfnTrain{} points (\BfnTrainCi) on the train months and \BfnHeld{} (\BfnHeldCi) held out, while
taking a bishop with a knight reads \NfbTrain{} (\NfbTrainCi) and \NfbHeld{} (\NfbHeldCi), across
0 on both. That fits a rule that holds only for keeping the bishop, and as well one that holds both
ways beside a lean towards the quiet move, which the mean cancels; the halves cannot tell the two
apart. The trade units side with the rule both ways. In each of them, on both periods, a trade
that takes a bishop fares better against the quiet move than one that takes a knight, and when
behind in development, the trade lifts by \BehindBishopHeld{} points held out where it takes a
bishop and loses by \BehindKnightHeld{} where it takes a knight, both intervals clear of 0 on both
periods.

Of trading by development's halves, trading when behind lifts by \BehindTrain{} points
(\BehindTrainCi) and \BehindHeld{} (\BehindHeldCi), and keeping the pieces on when ahead by
\AheadTrain{} (\AheadTrainCi) on the train months but \AheadHeld{} (\AheadHeldCi) held out. The
reference, a quiet move against an even trade of a piece with development level, reads
\LevelTrain{} (\LevelTrainCi) and \LevelHeld{} (\LevelHeldCi): with development level, neither the
quiet move nor the trade fares better than Maia-2's order says.

\UnitTable

\section{Where the two rules count}

\paragraph{Close calls.} A close call is a position where Maia-2 rates its second move at least half
as likely as its first: \CloseShareLow\% to \CloseShareHigh\% of a rule unit's held-out positions
(Table~\ref{tab:close}). There a tie-breaker matters most, since his level plays either move. In close calls, keeping the bishop where the first choice gives it for a knight
is sounder than that first choice by \BfnCloseBfTrain{} points (\BfnCloseBfTrainCi) on the train
months, and by \BfnCloseBfHeld{} (\BfnCloseBfHeldCi) held out, where the interval spans 0. A pawn
kept off the bishop's colour where the first choice puts one on it reads \ColCloseBfTrain{}
(\ColCloseBfTrainCi) and \ColCloseBfHeld{} (\ColCloseBfHeldCi), about as sound as that first
choice. Where the first choice is clear, it wins: there, following the rule against it costs
\BfnClearBfHeld{} points for keeping the bishop and \ColClearBfHeld{} for the bishop's colour held
out, below 0 on both periods.

A quiet move also fares better against an even trade in close calls when development is level, by
\LevelCloseTrain{} and \LevelCloseHeld{} points with both intervals above 0, and so does keeping the
tension on the train months, by \TenCloseTrain{} points, though not held out, at \TenCloseHeld. Part
of a quiet move's lift in a close call may be that lean, which this study does not separate. The
bishop's colour compares two quiet pawn moves, so no such lean touches it.

\CloseTable

\paragraph{Phase and band.} Keeping the bishop lifts most in the opening, by \BfnOpenHeld{} points
held out against \BfnMiddleHeld{} in the middlegame, where it still lifts on both periods, and more
in his band than in the band above, by \BfnHisHeld{} against \BfnAboveHeld{}
(Table~\ref{tab:splits}). Taking a bishop with a knight lifts in the opening, by \NfbOpenHeld, and
misleads in the middlegame, at \NfbMiddleHeld, both clear of 0 on both periods. The bishop's colour
lifts in the middlegame and the endgame, by \ColMiddleHeld{} and \ColEndHeld{} points held out,
above 0 on both periods, and not in the opening, at \ColOpenHeld{} with its interval across 0 on
both; it lifts in both bands.

\SplitTable

\section{Examples}

The rule draws six held-out examples a unit, at random among the positions where exactly one of
the two moves is sound. For keeping the bishop, the rule's move is the sound one in
\BfnExamplesF{} of its six and the other move in \BfnExamplesB; for the bishop's colour, in
\ColExamplesF{} and \ColExamplesB. Figure~\ref{fig:examples} shows three.

In the first, White, in \BfnFourBand, can give a bishop for a knight with \BfnFourB, which Maia-2
rates first, ahead of the quiet \BfnFourF. By the referee, White's win\% is \BfnFourFWin{} after
the quiet move and \BfnFourBWin{} after the trade; the referee's best is \BfnFourBest, and the game went
\BfnFourPlayed. In the second, Black, in \ColTwoBand, can put a pawn on the bishop's colour with
\ColTwoB, Maia-2's first choice, or keep it off with \ColTwoF, which was played and is the referee's
best: Black's win\% is \ColTwoFWin{} after it and \ColTwoBWin{} after \ColTwoB. The third is one
where the rule is wrong. Black, in \BfnSixBand, played the quiet \BfnSixF, Maia-2's first choice,
but giving the bishop for the knight with \BfnSixB{} is the referee's best: Black's win\% is
\BfnSixBWin{} after the trade and \BfnSixFWin{} after the quiet move.

\ExampleFigure

\section{What it means for the app}

Two of the coach's four rules carry information that his level's own order of moves lacks, and
they read best as he gives them: as the way to choose between two normal moves. A lesson can teach
them as tie-breakers. When two moves look about equally good and one gives the bishop for a knight,
play the other; when two pawn moves look about equally good, push the one that goes to the other
colour from the bishop. Neither is a reason to overrule a move that looks better: where his level's
first choice breaks the rule, that first choice is still sound more often than the rule's move, and
only in a close call is the rule's move about as sound. The bishop's colour counts in the
middlegame as in the endgame, where the app's catalogue already teaches it, and not in the opening.
Since chess PR 181 a strategy lesson teaches the two this way: it starts from checks, captures and
threats, plays a move that is clearly better whatever the rules say, and only then, between two
moves still in the balance, keeps the bishop (most of all in the opening) or pushes the pawn to the
other colour from it (in the middlegame and the endgame).

Keeping the tension and trading by development did not choose between his two natural moves, so a
lesson that teaches them explains a move, but should not offer them as the way to choose one.

\section{Caveats}

\begin{itemize}[nosep]
\item \textbf{Maia-2's order is the yardstick.} The balanced lift assumes the first choice's lead
  over the second is about the same in the two orders. A rule whose follow-first positions differ
  in kind from its break-first ones could lean either way, so each order's lift is reported, and
  the control shows the design sees a rule known to hold.
\item \textbf{Halves are not tests.} Only a rule's mean was tested. A half's lift can carry a lean
  towards quiet moves or towards trades, so the halves describe where a rule's lift sits, not which
  half holds.
\item \textbf{Many splits.} The splits are many, and some cells will read clear of 0 by chance; they
  describe and decide nothing.
\item \textbf{Not his games.} The positions are Lichess players' in his band and the band above;
  his own games are too few for a test like this.
\item \textbf{Two candidates.} The coach's student names a few normal moves. The test takes the
  two Maia-2 rates most likely at the mover's band, and only where one follows the rule and the
  other breaks it.
\item \textbf{Quiet by board code.} No mate in one and no winning capture for either side is the
  app's reading of a quiet moment. A deeper tactic, a fork two moves away, still counts as quiet.
\item \textbf{Even by the capture war.} A trade the app's capture war calls even can still lose to a
  tactic it does not see; that is the referee's to find, and part of what a rule's move is judged on.
\item \textbf{Development by the starting squares.} The app has no reading of development, so the
  count of undeveloped minor pieces and the king on its square is this study's. A rook not yet
  connected or a queen out early does not count.
\item \textbf{One move for a plan.} A rule is tested on the next move, not as a plan over several,
  and a sound move can still be the second-best plan.
\item \textbf{Sound within a window.} A move is sound within \WindowWp{} win\% of the best; a rule
  whose gain is smaller than that, or comes later than the referee's search sees, reads as none.
\end{itemize}

\end{document}
